\documentclass[preview]{standalone}

\usepackage{amsmath}
\usepackage{amssymb}
\usepackage{stellar}
\usepackage{bettelini}
\usepackage{definitions}
\usepackage{xfrac}

\begin{document}

\id{series-functions-exercises}
\genpage

\section{Exercises}

\begin{snippetexercise}{series-of-functions-ex1}{}
    Consider the series of \function[functions]
    \[
        \sum_{n=0}^\infty x^n (1-x)^n
    \]
    Determine the \set \(S\subseteq [0,1]\) in which the series converges pointwise.
    Determine whether the sum of the series is a continuous function on \([0,1]\).
\end{snippetexercise}

\begin{snippetsolution}{series-of-functions-ex1-sol}{}
    Let \[ f_n(x) = x^n(1-x)^n = ((x(1-x)))^n, \quad x \in [0,1] \]
    On \([0,1]\), the term \(0 \leq x(1-x) \leq \sfrac14\) since the maximum is attained at \(x=\sfrac12\).
    Thus, \[
        \sum_{n=0}^\infty x^n (1-x)^n \leq \sum_{n=0}^\infty \left(\frac14\right)^{n}
    \]
    which converges. The pointwise convergence set is \(S = [0,1]\).
    The sum is given by
    \[
        \sum_{n=0}^\infty x^n (1-x)^n = \frac{1}{
            1 - x(1-x)
        } = \frac{1}{1 - x + x^2}
    \]
    which is continuous on \(S\).
    We would have otherwise noted that the series converges uniformly, and the sum of continuous terms is continuous.
\end{snippetsolution}

\begin{snippetexercise}{series-of-functions-ex2}{}
    Determine the radius of converges of following power series:
    \[
        \sum_{n=0}^\infty \frac{n^2}{3^n}x^n
    \]
    Study the convergence on the boundary.
\end{snippetexercise}

\begin{snippetsolution}{series-of-functions-ex2-sol}{}
    The radius of convergence is given by
    \begin{align*}
        R = \left(\limsup_n \left|\frac{n^2}{3^n}\right|\right)^{-1}
        &= \left(\frac13 \cdot \limsup_n n^{\sfrac{2}{n}}\right)^{-1} = 3
    \end{align*}
    Thus the power series converges absolutely for \(|x|<3\).
    For \(x=3\) we have
    \[
        \sum_{n=0}^\infty \frac{n^2}{3^n}3^n = \sum_{n=0}^\infty n^2
    \]
    which diverges. Likewise, for \(x=-3\) we have
    \[
        \sum_{n=0}^\infty \frac{n^2}{3^n}{(-3)}^n = \sum_{n=0}^\infty (-1)^n n^2
    \]
    but the general term does not tend to zero, thus it also diverges.
\end{snippetsolution}

\begin{snippetexercise}{series-of-functions-ex3}{}
    Determine the radius of converges of following power series:
    \[
        \sum_{n=1}^\infty \left(
            1 + \frac1n
        \right)^{n^2}x^n
    \]
\end{snippetexercise}

\begin{snippetsolution}{series-of-functions-ex3-sol}{}
    The radius of convergence is given by
    \[
        R = \left(
            \limsup_n \left(1 + \frac1n\right)^n
        \right)^{-1} = \frac{1}{e}
    \]
\end{snippetsolution}

\begin{snippetexercise}{series-of-functions-ex4}{}
    Determine the radius of converges of following power series:
    \[
        \sum_{n=1}^\infty \frac{
            (nx)^n
        }{n!}
    \]
\end{snippetexercise}

\begin{snippetsolution}{series-of-functions-ex4-sol}{}
    The radius of convergence is given by
    \begin{align*}
        \frac1R &= \lim_n \left|\frac{a_{n+1}}{a_n}\right|
        = \lim_n \frac{(n+1)^{n+1}}{(n+1)!} \cdot \frac{n!}{n^n} \\
        &= \lim_n \frac{(n+1)^n}{n^n} = \lim_n \left(1 + \frac1n\right)^n = e
    \end{align*}
    Thus, \(R = \sfrac1e\).
\end{snippetsolution}

\begin{snippetexercise}{series-of-functions-ex5}{}
    Consider the series of \function[functions]
    \[
        f(x) = \sum_{n=3}^\infty x^{-\log n}
    \]
    \begin{enumerate}
        \item Determine the natural domain \(D\) of \(f\), and specify
        in which subintervals of \(D\) the series converges uniformely.
        \item Trace a qualitative graph of \(f\). In particular, show that
        \(f \in \mathcal C^1(D)\). The study of the second derivative is not necessary.
        \item Determine whether \(f\) is integrable over \(D\), even as an improper integral.
    \end{enumerate}
\end{snippetexercise}

\begin{snippetsolution}{series-of-functions-ex5-sol}{}
    We have for \(x>0\)
    \begin{align*}
        x^{-\log n} &= e^{\log(x^{-\log n})} = e^{-(\log n)(\log x)} = n^{-\log(x)}
    \end{align*}
    This is a \(p\)-series with \(p = \log x\) which converges \ifandonlyif
    \(\log x < 1\), meaning \(x > e\). The natural domain is \(D = (e, +\infty)\).
    TODO.    
\end{snippetsolution}

\begin{snippetexercise}{series-of-functions-ex6}{}
    Determine the radius of converges of following power series:
    \[
        \sum_{n=1}^\infty \frac{
            n!
        }{
            (2n)!
        }x^n
    \]
\end{snippetexercise}

\begin{snippetsolution}{series-of-functions-ex6-sol}{}
    The radius of convergence is given by
    \begin{align*}
        \frac1R &= \lim_n \left[\frac{a_n}{a_{n+1}}\right]
        = \lim_n \frac{(n+1)!}{(2n+2)!} \cdot \frac{(2n)!}{n!} = \\ 
        &= \lim_n
        \frac{(n+1)n!}{(2n+2)(2n+1)(2n)!} \cdot \frac{(2n)!}{n!}
        = \lim_n \frac{1}{2(2n+1)} = 0
    \end{align*}
    Thus, \(R=\infty\).
\end{snippetsolution}

\begin{snippetexercise}{series-of-functions-ex7}{}
    Determine the radius of converges of following power series:
    \[
        \sum_{n=1}^\infty 2^nx^{2n}
    \]
\end{snippetexercise}

\begin{snippetsolution}{series-of-functions-ex7-sol}{}
    We have
    \[
        \sum_{n=1}^\infty 2^nx^{2n}
        = \sum_{n=1}^\infty (2x^2)^n
    \]
    which is a geometric series \(\sum r^n\) which converges \ifandonlyif
    \(|2x^2| < 1\), meaning \(|x| < \sfrac{1}{\sqrt 2}\).
    Thus, \[
        R = \frac{1}{\sqrt{2}}
    \]
\end{snippetsolution}

\begin{snippetexercise}{series-of-functions-ex8}{}
    Determine the radius of converges of following power series:
    \[
        \sum_{n=1}^\infty \frac{n}{2^{2n}}x^n
    \]
\end{snippetexercise}

\begin{snippetsolution}{series-of-functions-ex8-sol}{}
    The radius is given by
    \[
        \frac1R = \lim_n \frac{n+1}{2^{2n+2}} \cdot \frac{2^{2n}}{n} = \lim_n \frac{n+1}{4n} = \frac14
    \]
    Thus, \(R = 4\).
\end{snippetsolution}

\begin{snippetexercise}{series-of-functions-ex9}{}
    Determine the radius of converges of following power series:
    \[
        \sum_{n=1}^\infty n!x^{n^2}
    \]
    \emph{Hint:} use Stirling's formula.
\end{snippetexercise}

\begin{snippetsolution}{series-of-functions-ex9-sol}{}
    Consider
    \[
        \sum_{n=1}^\infty n!x^{n^2}
        = \sum_{m=0}^\infty a_m x^m, \quad a_m = \begin{cases}
            n! & m = n^2 \text { for some } n \\
            0 & \text{otherwise}
        \end{cases} 
    \]
    The radius of convergence is given by
    \begin{align*}
        \frac1R &= \limsup_m |a_m|^{\sfrac1m}
        = \limsup_n (n!)^{\sfrac1{n^2}}
        = \limsup_n \left(
            \sqrt{2\pi n}\left(\frac ne\right)^n
        \right)^{\sfrac1{n^2}} \\
        &= \limsup_n \frac{
            (2\pi n)^{\sfrac{1}{2n^2}} n^{\sfrac1n}
        }{e^{\sfrac1n}} \to 1
    \end{align*}
    Thus, \(R = 1\).
\end{snippetsolution}

\begin{snippetexercise}{series-of-functions-ex10}{}
    Show that for every \(z \in \complexnumbers\), we have
    \[
        |e^z - 1| \leq e^{|z|} - 1 \leq |z| e^{|z|}
    \]
\end{snippetexercise}

\begin{snippetsolution}{series-of-functions-ex10-sol}{}
    \todo
\end{snippetsolution}

\begin{snippetexercise}{series-of-functions-ex11}{}
    Show that the series of \function[functions]
    \[
        \sum_{n=0}^\infty e^{nx}
    \]
    converges for every \(x<0\) and calculate its sum.
    Determine for which \set[sets] the series converges uniformely.
\end{snippetexercise}

\begin{snippetsolution}{series-of-functions-ex11-sol}{}
    \todo
\end{snippetsolution}

\includesnpt{integral-sine-function-definition}
\includesnpt{integral-sine-function-maclaurin-series}

\begin{snippetexercise}{series-of-functions-ex12}{}
    Study the convergence of the series of \function[functions]
    \[
        \sum_{n=3}^\infty \frac{
            {(\log n)}^{nx}
        }{n!}
    \]
\end{snippetexercise}

\begin{snippetsolution}{series-of-functions-ex12-sol}{}
    \todo
\end{snippetsolution}

\begin{snippetexercise}{series-of-functions-ex13}{}
    Let \(\varphi \in C_b(\realnumbers)\), \(f_0 = \varphi\), \(f_{n+1}(x) = \integral[0][x][f_n(t)][t]\).
    Study the convergence of
    \[
        \sum_{n=0}^\infty f_n(x)
    \]
    Determine the sum of the series.
\end{snippetexercise}

\begin{snippetsolution}{series-of-functions-ex14-sol}{}
    \todo
\end{snippetsolution}

\begin{snippetexercise}{series-of-functions-ex14}{}
    Determine the natural domain \(D\) of the \function
    \[
        f(x) = \sum_{n=0}^\infty \frac{e^{-nx}}{1 + n^2x^2}
    \]
    Determine the subintervals of \(D\)
    in which the series converges uniformly
    and whether \(f \in \mathcal C^1(D)\).
\end{snippetexercise}

\begin{snippetsolution}{series-of-functions-ex14-sol}{}
    \todo
\end{snippetsolution}

\end{document}